\section{Soft sync 的耦合动力学}

若 anchor 使用指数滑动平均（exponential moving average，EMA），其概率空间理想化更新率记为 $\zeta$，采样提议分布记为 $\beta$。令
\begin{equation*}
 q_{k+1}=(1-\zeta)q_k+\zeta\widehat p_{k,N_k}
\end{equation*}
表示概率空间中的近似 soft sync。对 PMD-Part、BPO 和 MDPO，精确 inner 目标是
\begin{equation*}
 \widehat p_{k,N_k}\propto q_k^{1-\eta}t^\eta.
\end{equation*}
在固定 prefix 上令 $q=q_k$、$z=\log(q/t)$，假设 $\max_a|z_a|\le\varepsilon$，$\eta$ 固定且 inner 精确求解。归一化 $\E_te^z=1$ 的 Taylor 展开给出
\begin{equation*}
\E_tz=-\frac12\E_tz^2+O(\varepsilon^3),\qquad
\kappa=\E_t[e^zz]=\frac12\E_tz^2+O(\varepsilon^3).
\end{equation*}
因此 OPD 的混合配分函数满足
\begin{equation*}
\log Z_\eta=\log\E_te^{(1-\eta)z}
=-\frac{\eta(1-\eta)}2\E_tz^2+O(\varepsilon^3).
\end{equation*}
对于 Gibbs 更新，$\log(p^\star/t)=(1-\eta)z-\log Z_\eta$，一阶项为 $(1-\eta)z$。

PMD-Mean 的精确更新使用 Lambert-$W$ target。由式~\eqref{dyn:eq:mean-opd-renyi-bounds}，
\begin{equation*}
0\le\lambda\le\eta\kappa+\log Z_\eta
=\frac{\eta^2}2\E_tz^2+O(\varepsilon^3)=O(\varepsilon^2).
\end{equation*}
又因 $G=\eta(-z+\kappa)=O(\varepsilon)$，有
$W_0(\lambda e^G)=O(\varepsilon^2)$。故 Mean 也满足
\begin{equation*}
\log(p^\star/t)=z+G-W_0(\lambda e^G)
=(1-\eta)z+O(\varepsilon^2).
\end{equation*}
两类精确 inner 的概率空间混合 $q_{\mathrm{new}}=(1-\zeta)q+\zeta p^\star$，其中 $\zeta\in(0,1]$，均满足
\begin{align*}
q_{\mathrm{new}}/t
&=(1-\zeta)e^z+\zeta e^{(1-\eta)z+O(\varepsilon^2)}\\
&=1+(1-\zeta\eta)z+O(\varepsilon^2),\\
\log(q_{\mathrm{new}}/t)
&=(1-\zeta\eta)z+O(\varepsilon^2).
\end{align*}
这是 teacher 附近的局部展开。Mean 的精确全局递推仍由 Lambert-$W$ 决定。较小的 $\zeta$ 减缓 anchor 漂移和局部 teacher tracking；有限 inner 或参数 EMA 还需要加入相应误差项。

\begin{proposition}[概率 soft sync 的局部稳定条件]
\label{dyn:soft-local}
固定有限词表、严格正 teacher、$\eta>0$ 与 $\zeta\in(0,1]$。精确 Gibbs 或 Mean inner 后执行上述概率混合时，teacher 中心坐标满足
\begin{equation}
y_{\mathrm{new}}=(1-\zeta\eta)y+O(\|y\|_\infty^2).
\end{equation}
当 $0<\zeta\eta<2$ 时 teacher 局部渐近稳定；当 $\zeta\eta>2$ 时局部不稳定。边界 $\zeta\eta=2$ 需要高阶分析。本命题针对概率混合，不直接适用于参数 EMA。
\end{proposition}
\begin{proof}
在中心坐标中 $q=t e^y/\E_te^y$ 且 $\E_ty=0$，因此 $q/t=1+y+O(\|y\|_\infty^2)$。Gibbs 的 teacher 中心输出是 $(1-\eta)y$；Mean 的三阶展开给出相同一阶项。因此两种 target 均满足 $p^\star/t=1+(1-\eta)y+O(\|y\|_\infty^2)$。概率混合给出
\begin{equation*}
\frac{q_{\mathrm{new}}}{t}
=(1-\zeta)\frac qt+\zeta\frac{p^\star}t
=1+(1-\zeta\eta)y+O(\|y\|_\infty^2).
\end{equation*}
取对数并应用 $C_t$，公共常数消去，得到所述展开。余项在固定 teacher 的有限维邻域中一致有界。若 $|1-\zeta\eta|<1$，取充分小邻域使余项范数不超过 $(1-|1-\zeta\eta|)\|y\|_\infty/2$，则每步范数乘子严格小于一，邻域不变且误差几何趋零。若 $|1-\zeta\eta|>1$，取余项不超过 $(|1-\zeta\eta|-1)\|y\|_\infty/2$，更新范数严格放大，固定点局部不稳定。等于一时这些估计无法给出严格收缩或放大，所以边界不由线性项决定。
\end{proof}

\section{使用旧 anchor target 的附加误差}

若实际优化的是旧 anchor $q_j$ 对应的 Gibbs target，而当前 outer anchor 已为 $q_k$，则对任意 $\eta>0$，归一化常数消去后
\begin{equation*}
 \operatorname{osc}_v
 \log\frac{p_j^{\mathrm G}(v)}
          {p_k^{\mathrm G}(v)}
 =|1-\eta|\operatorname{osc}_v\log\frac{q_j(v)}{q_k(v)}.
\end{equation*}
因此对当前 target 的总跨度误差，至多是相对旧 target 的优化跨度误差加上右侧的 stale-anchor 项。teacher 与 $\eta$ 在这两个 target 间固定是该等式的前提。这个 target 漂移与估计量章节中旧 rollout 的采样测度偏差是两个独立问题，分别进入 outer 递推和梯度估计。

\section{访问测度固定的含义}

本章所有递推都是逐 history 的 token distribution 递推。真实 LLM 中 history 集合由 $d_k$ 产生；只要 inner 内固定该集合，以上结论适用于 conditional token target。outer rollout 更新后，$d_{k+1}$ 改变，必须额外记录 coverage drift：
\begin{equation*}
 \varepsilon_k^{\mathrm{visit}}=\TV(d_{k+1},d_k).
\end{equation*}
若 token loss 在所有 history 上有上界 $C$，则访问测度改变使目标值的期望变化至多 $2C\varepsilon_k^{\mathrm{visit}}$；若要控制梯度的期望变化，需另行假设逐 prefix 梯度向量的范数不超过 $C_g$，此时界为 $2C_g\varepsilon_k^{\mathrm{visit}}$。因而“固定测度下求解 token 子问题”和“outer 重新 rollout”是两个严格分开的步骤。
